% ============================================================
% AIR SOURCE ABSTRACTION JUSTIFICATION — Team 0109D — 04.4_Verification
% ============================================================
\documentclass[12pt, letterpaper]{article}
\usepackage[margin=0.5in]{geometry}
\usepackage[T1]{fontenc}
\usepackage[scaled=1]{helvet}
\renewcommand{\familydefault}{\sfdefault}
\usepackage[utf8]{inputenc}
\usepackage{booktabs}
\usepackage{amsmath}
\usepackage{enumitem}
\usepackage{hyperref}
\usepackage{xcolor}
\usepackage{titlesec}

\titleformat{\section}{\large\bfseries}{\thesection}{0.6em}{}
\titlespacing*{\section}{0pt}{8pt}{3pt}
\setlist{nosep, leftmargin=1.2em, itemsep=2pt}
\hypersetup{colorlinks=true, urlcolor=blue!60!black}
\setlength{\parskip}{4pt}
\setlength{\parindent}{0pt}

\begin{document}

\begin{center}
{\Large\textbf{Air Source Abstraction Justification}}\\[2pt]
{\normalsize Team 0109D \quad$\vert$\quad ESC204 --- Praxis III \quad$\vert$\quad Winter 2026}
\end{center}
\noindent\rule{\textwidth}{0.6pt}\vspace{6pt}

\section{Design Concept vs.\ Prototype Abstraction}

The final Design Concept specifies a dedicated \textbf{air compressor} as the pressure source, delivering airflow through ducting to the ridge-mounted nozzle carriage. For the Phase~3 prototype, this was abstracted to a \textbf{Dyson hair dryer} (max airflow setting, cool mode). This section justifies why this abstraction is valid and does not compromise the verification results.

\section{Why Abstract the Air Source?}

The prototype's purpose is to verify two independent hypotheses:

\begin{enumerate}
    \item \textbf{Mechanical translation:} Can a stepper-driven belt carriage reliably traverse a ridge track? --- This is independent of the air source entirely.
    \item \textbf{Nozzle fluid dynamics:} Does a planar slit geometry outperform a circular geometry for surface clearing? --- This depends on \textit{flow characteristics at the nozzle exit}, not on what generates the flow upstream.
\end{enumerate}

The air compressor is a \textit{commodity component selection problem} (choosing a compressor that meets CFM/PSI specs), not a \textit{design verification problem}. Procuring, plumbing, and integrating a compressor within the prototype budget and timeline would consume resources without testing any novel engineering hypothesis. The Dyson provides continuous, high-volume airflow that allows us to isolate and test the nozzle geometry --- the actual design variable.

\section{Why the Dyson Is a Valid Stand-In}

\begin{table}[h]
\centering\small
\begin{tabular}{@{}lll@{}}
\toprule
\textbf{Property} & \textbf{Final Compressor} & \textbf{Dyson (Prototype)} \\
\midrule
Flow type & Continuous & Continuous \\
Volume (CFM) & High (design spec TBD) & High ($\sim$13 L/s) \\
Pressure delivery & Steady-state & Steady-state \\
Force output & Design-dependent & Measured: $\sim$1.96\,N at 12\,cm \\
Controllability & On/off via relay & On/off manual (sufficient for test) \\
\bottomrule
\end{tabular}
\end{table}

The critical parameter is \textbf{continuous, steady-state volumetric flow} through the nozzle. Both the Dyson and a compressor deliver this. The Dyson was explicitly chosen over alternatives like a ball pump (pulsed, low-volume, unsuitable for boundary-layer shear) or a shop vacuum (variable pressure, difficult to control).

\section{What This Abstraction Validates and Does Not Validate}

\textbf{Validated by prototype testing:}
\begin{itemize}
    \item Nozzle geometry comparison (planar slit vs.\ circular) at equivalent input flow
    \item 30$^\circ$ angle of attack optimization (Coanda effect, $F_x$/$F_y$ balance)
    \item Wet debris stiction clearing capability at distance
    \item Force conservation across nozzle geometries (mass-flow conservation)
\end{itemize}

\textbf{Deferred to next phase (requires compressor selection):}
\begin{itemize}
    \item Final CFM/PSI specification for the compressor
    \item Noise compliance (O-3.1-STRU-1: $<$70\,dBA at 3\,m) --- the Dyson exceeds this threshold, but it is not the production air source
    \item Ducting pressure losses from compressor to ridge
    \item Weatherproofing of the air supply path
\end{itemize}

\section{Precedent}

This abstraction follows standard prototyping practice: isolate the design variable (nozzle geometry) from commodity components (air source) to focus verification effort on the novel engineering contribution. The team consensus to pivot from the originally planned ball pump to the Dyson was made early in Phase~3 to ensure all verification tests could use a continuous, high-volume source representative of the final system's operating regime.

\end{document}
